% !PRINTED 109
\weberpage{109}

For if we put
\begin{equation}
 f(x)=x^n-a,
\tag{1}
\end{equation}
then by § 31, VII, $f(x)$ is positive for sufficiently large positive $x$, and negative for $x=0$; hence there is a positive value of $x$, denoted by $\sqrt[n]{a}$, for which $f(x)$ vanishes. That there can be only one such value follows from the fact that $x^n$, as long as $x$ remains positive, increases continually with $x$.

If $n$ is odd, then also for negative $a$ there exists one and only one negative root of (1), which is proved in the same way.

\medskip
\noindent X. \gesperrt{An equation of odd degree with real coefficients always has at least one real root.}
\medskip

For by theorem VII, $f(x)$, if the degree is odd, can become both positive and negative; hence, by VIII, $f(x)=0$ has a root.

Finally we prove:

\medskip
\noindent XI. \gesperrt{A pure equation always has a root.}
\medskip

By a pure equation we understand an equation of the form
\begin{equation}
 x^n-a=0,
\tag{2}
\end{equation}
where $a$ is an arbitrary complex magnitude. That this equation always has a root for real positive $a$ has already been proved above; and for $a=0$ it is plainly satisfied by $x=0$.

The general proof can be divided into two parts; it is sufficient to assume first $n=2$, and then $n$ odd. For if the magnitudes $\sqrt[n]{a}$ and $\sqrt[2]{a}$, or $\sqrt{a}$, exist for every $a$, then the existence of
\[
 \sqrt[2n]{a}=\sqrt[2]{\sqrt[n]{a}},\qquad \sqrt[4n]{a}=\sqrt[2]{\sqrt[2n]{a}}\ldots,
\]
and hence of $\sqrt[n]{a}$ for every arbitrary $n$, is secured. It would even suffice to prove the existence of $\sqrt[n]{a}$ for the case in which $n$ is a prime number. But from this no particular advantage is at first to be drawn.

We denote the complex magnitude $a$ by $b+ci$ and put, since $x$ will in general also be complex,
\[
 x=y+iz.
\]
